\begin{theorem}[G-24: The exact eventual rectangular range through $m$ rows]\label{G-24}
Let $n>m\ge17$ and $4\le r\le m$. At every sufficiently
large admissible coordinate degree $d$ (meaning $r\mid nd$),
the rectangular multiplicity gap
\[
 m_{(nd/r)^r}(\ell^{n-m}\operatorname{per}_m)
       >m_{(nd/r)^r}(\det_n)
\]
holds if and only if
\begin{equation}\label{G24:range}
 m<n<\frac{rm}{r-1}.
\end{equation}
Inside this range, G-23 supplies the explicit threshold.
Above the boundary, the permanent rectangular multiplicity
is zero in every positive admissible degree and the determinant
multiplicity is positive in every sufficiently large admissible
degree. At the equality boundary, the determinant multiplicity
is strictly larger in every sufficiently large degree.
No assertion excluding all earlier boundary gaps is made for
$r\ge5$; G-12 supplies that stronger weak inequality for $r=4$.

For an explicit equality-boundary threshold, write
$a=n-m$, so that $n=ra$, $m=(r-1)a$, and define
\[
 \begin{split}
 u&=(r-2)n^2-r^2+3,\\
 v&=\begin{cases}
       \binom{m+2}{2}-8,&r=4,\\
       (r-1)m^2-2m-r^2+2r+2,&r\ge5.
      \end{cases}
 \end{split}
\]
Let
\[
 b=\binom{n+r-1}{r-1},\qquad
 \beta=\left\lceil\frac{3r^2b}{8}\right\rceil n^{2(r^2-1)},
 \qquad T=\binom{b+\beta}{b}-1,\qquad E=rT(\beta!).
\]
Then $u>v\ge2$, and determinant dominance holds for every
\begin{equation}\label{G24:threshold}
 d\ge E\bigl(4v(nE)^{v-1}+1\bigr).
\end{equation}
\end{theorem}

\begin{proof}
The strict interior is the rectangular consequence of G-23.
For the other cases, set $a=n-m$ and use $r$ variables
$x_0,\ldots,x_{r-1}$. A nonzero padding form can be taken
to $x_0$ by a special-linear change of variables. The
one-parameter subgroup
\[
 \lambda(t)=\operatorname{diag}(t^{r-1},t^{-1},\ldots,t^{-1})
\]
gives weight $(r-1)a-m+rj=ra-n+rj$ to a monomial of
$x_0^a q$ coming from the $x_0^j$ part of $q$.
If $(r-1)n>rm$, all these weights are positive. The form
therefore tends to zero and every positive-degree invariant
vanishes. This proves the zero assertion above the boundary.
For eventual determinant occurrence, take the common free
point of R-28 with smaller size $n-1$ and padding one.
For the minimal admissible direction use
$s_0=r/\gcd(n,r)$ and $\nu=(ns_0/r)^r$.
Its strict condition is $ns_0/r>s_0$, since $n>r$.
R-28 consequently gives positive determinant multiplicity
at every sufficiently large integer stretch, that is, every
sufficiently large admissible coordinate degree.

\paragraph{1. Invariant functions inject into one fewer variables at equality.}
Now suppose $(r-1)n=rm$. Then $n=ra$ and $m=(r-1)a$.
The same limit is
\begin{equation}\label{G24:limit}
 \lim_{t\to0}\lambda(t)x_0^a q
                =x_0^a q(0,x_1,\ldots,x_{r-1}).
\end{equation}
Put $q_0=r-1$ for the number of variables after restriction,
and write
\[
 A_P=\mathbb C[\{L^a q:q\in\mathcal P_{m,r}\}^{\rm cl}]^{\mathrm{SL}_r},
 \qquad C=\mathbb C[\mathcal P_{m,q_0}]^{\mathrm{SL}_{q_0}}.
\]
Restriction along $q\mapsto x_0^a q$ gives a graded algebra
homomorphism
\begin{equation}\label{G24:injection}
 \rho:A_P\longrightarrow C.
\end{equation}
It is well defined because the indicated padded forms belong
to the original cone, and the subgroup
$\operatorname{diag}(1,\mathrm{SL}_{q_0})$ fixes $x_0$.
Its grading is unchanged: coefficients of $x_0^a q$ depend
linearly on coefficients of $q$.

To prove injectivity, let an invariant function $h$ restrict
to zero. For any original padded form, take its nonzero
padding line to $x_0$. The function $h$ is constant along
$\lambda(t)$ and hence, by regularity, at its limit
\eqref{G24:limit}. Restricting a permanent pencil to
$x_0=0$ puts the residual form in $\mathcal P_{m,q_0}$,
including its closure points. Thus the value at the limit
is zero. The dense padded image has zero value everywhere,
so $h=0$ on its closure. The same statement holds in degree
zero. This proves \eqref{G24:injection} without claiming
that every lower-rank invariant lifts.

Let $A_D=\mathbb C[\mathcal D_{n,r}]^{\mathrm{SL}_r}$.
For $r\ge5$ put $B=C$. For $r=4$ use instead
\[
 B=\mathbb C[\operatorname{Sym}^m\mathbb C^3]^{\mathrm{SL}_3}.
\]
In this latter case exactness of invariants gives a graded
surjection $B\to C$, not an asserted injection into the
full ternary space. Consequently in every case
\begin{equation}\label{G24:hilbertupper}
 \dim(A_P)_d\le\dim B_d.
\end{equation}

\paragraph{2. The two comparison rings have known dimensions and no period.}
The equality hypotheses and $r\le m$ imply $a\ge2$.
Apply R-28's common-free-point statement, which allows
row count three, first with degree $n$, size $n-1$, row
count $r$ and rectangular direction $(a^r)$ in degree one.
Apply it a second time with degree $m$, size $m-1$, row
count $q_0$ and direction $(a^{q_0})$.
Both applications satisfy the strict inequality $a>1$ and
the size condition. They supply a stable point with no
nontrivial scalar-orbit symmetry in $A_D$ and in the
permanent cone underlying $C$. The latter is also such a
point in the full ambient ternary representation when
$r=4$.

It follows that the stable quotients have dimensions
\[
 \dim A_D=(r-2)n^2+2-(r^2-1)=u,
 \qquad \dim B=v.
\]
For $r\ge5$, use G-21 in $q_0\ge4$ variables; for $r=4$
the whole ternary coefficient space has dimension
$\binom{m+2}{2}$ and the stable orbit dimension is eight.
The absence of scalar-orbit symmetry and invariant
separation give gcd one for the nonzero invariant degrees,
just as in R-28. We use this separately for the two rings;
no common invariant in their different ambient spaces is
being assumed.

The dimension inequality is strict. For $r\ge5$ it reads
\[
 u-v=\frac{r^2-3r+1}{(r-1)^2}m^2+2m-2r+1>0,
\]
since $m\ge r$ and $r^2-3r+1>0$.
For $r=4$ it is
\[
 u-v=\frac{55m^2-27m-108}{18}>0\qquad(m\ge17).
\]
The displayed dimensions are at least two in the stated range.

\paragraph{3. An explicit comparison at every integer degree.}
Derksen's general bound, specialized as in G-23 with $s=1$
and $\eta=0$, gives the displayed $\beta$ for the degree-$n$
ambient representation in $r$ variables. It also bounds the
generators for the degree-$m$ representation in $q_0<r$
variables: its vector-space dimension, action degree and
group-dimension exponent are all no larger. Exactness of
invariants passes these bounds to $A_D$ and to $B$.
The chosen $T$ bounds both generator counts. Hence the
$E$th Veroneses of both rings are standard graded by the
same generator-grouping argument as G-23.

The gcd-one argument and the bound $\beta$ give nonzero
elements in every degree at least $\beta^2$ in each ring
separately. For $d=Ek+t$, $0\le t<E$, multiplication by
elements of degrees $E+t$ and $2E-t$ therefore gives the
individual sandwiches
\[
 \dim A_{E(k-1)}\le\dim A_d\le\dim A_{E(k+2)},
 \qquad A=A_D\text{ or }B.
\]
The projective degree of $B^{(E)}$ is at most
$(nE)^{v-1}$. If $r\ge5$, compose its invariant-coordinate
map with the permanent-pencil parametrization of degree
$m<n$ and use G-10. If $r=4$, use the identity parametrization
of the full ternary coefficient space, of degree one, and
the same larger bound. No dominance of the ternary
permanent parametrization is needed here.

Noether normalization on $A_D^{(E)}$, whose dimension is
at least $v+1$, and the degree/Hilbert bound on $B^{(E)}$
now give
\[
 \dim(A_D)_{E(k-1)}\ge\binom{k+v-1}{v},\qquad
 \dim B_{E(k+2)}\le(nE)^{v-1}\binom{k+v+1}{v-1}.
\]
The same exact ratio as in G-23 is at least $k/(4v)$ for
$k\ge v+1$. For $k\ge4v(nE)^{v-1}+1$, this proves
$\dim(A_D)_d>\dim B_d\ge\dim(A_P)_d$ in every residue.
All degrees are admissible at this boundary because $r\mid n$.
The invariant multiplicities are the claimed rectangles,
by the same single-label extraction and inheritance as
G-11/G-22. This proves \eqref{G24:threshold} and the theorem.
\end{proof}

\paragraph{Exact scope.}
The equality-boundary argument is an injection of actual
invariant coordinate rings followed by a Hilbert-dimension
comparison. It neither identifies those rings nor replaces
orbit-closure multiplicities by ambient plethysm coefficients.
The theorem closes the eventual rectangular boundary through
$m$ rows in the stated size range. It leaves nonrectangular
boundary faces, earlier finite-degree gaps at $r\ge5$,
more than $m$ rows, and efficient positive extraction open.

